\section{Negative branch and the scalar side of the problem}
\label{sec:negative}

The negative branch in Eq.~\eqref{eq:roots} is fixed by the same secular equation as the positive branch.  It should therefore be retained in the analytical structure even if the W/Z ratio uses only \(x_+\).  The root identities imply
\begin{equation}
  x_-(J)=-J-x_+(J)=-\frac{J}{x_+(J)}.
  \label{eq:negative-from-positive}
\end{equation}
Any physical use of the positive branch automatically carries a partner branch.

The Higgs potential can be written as
\begin{equation}
  V(H)=m^2 H^\dagger H+\lambda (H^\dagger H)^2,
  \qquad m^2<0
  \label{eq:higgs-potential}
\end{equation}
in the broken phase.  The sign of the quadratic term makes the negative branch tempting as an order-parameter datum.  A valid identification must be gauge-invariant and scheme-controlled.  Possible targets include \(m^2\), \(\lambda v^2\), the Higgs pole mass, or an eigenvalue in a compactification problem.

The branch values associated with \(J=3/4\) and \(J=2\) generate definite dimensionless numbers.  Multiplying them by a scale fitted to the vector masses gives scales near the electroweak/Higgs range.  This observation is a prompt for analysis.  The manuscript should not promote it to a prediction until a scalar-sector map is derived.

A useful criterion is the following.  If the same two-channel determinant produces both the vector pole ratio and the scalar instability, then the determinant should appear in an effective action, boundary condition, or self-energy equation.  The next calculation is to write such a determinant explicitly and identify which field or mode occupies the second channel.

\paragraph*{Section status.}  The negative branch algebra is exact.  Its Higgs interpretation is conjectural and requires a dynamical map.
